\documentclass[a4paper,11pt]{article}

\newcommand{\TPnumero}{Project 3}
\newcommand{\macrosPath}{../../preamble/macros.tex}
% =========================================================
% Tutorial / lab preamble — Time Series Analysis module
% article class + fancyhdr + tcolorbox + listings (English)
% Author : Aymen Ben Brik (aymen.benbrik@esprit.tn)
% =========================================================

\usepackage[utf8]{inputenc}
\usepackage[T1]{fontenc}
\usepackage[english]{babel}
\usepackage{lmodern}
\usepackage{amsmath, amssymb, amsthm, mathtools, bm}
\usepackage{graphicx}
\usepackage{booktabs}
\usepackage{array}
\usepackage{multirow}
\usepackage{colortbl}
\usepackage{xcolor}
\usepackage{tikz}
\usetikzlibrary{positioning, arrows.meta, calc, shapes, fit, backgrounds, matrix}
\usepackage{listings}
\usepackage[most]{tcolorbox}
\usepackage{geometry}
\geometry{a4paper, margin=2cm, headheight=15pt}
\usepackage{fancyhdr}
\usepackage[hidelinks]{hyperref}
\usepackage{enumitem}

% --- Colours ---
\definecolor{espritBlue}{RGB}{30, 60, 110}
\definecolor{espritAccent}{RGB}{200, 80, 30}
\definecolor{boxEx}{RGB}{30, 60, 110}
\definecolor{boxQ}{RGB}{180, 120, 0}
\definecolor{boxC}{RGB}{30, 100, 60}
\definecolor{boxAtt}{RGB}{200, 80, 30}

% --- Listings ---
\definecolor{codebg}{RGB}{248,248,248}
\definecolor{codeframe}{RGB}{220,220,220}
\definecolor{keyword}{RGB}{0,82,155}
\definecolor{comment}{RGB}{88,110,117}
\definecolor{string}{RGB}{133,153,0}

\lstdefinestyle{pythonstyle}{
  language=Python,
  basicstyle=\ttfamily\small,
  keywordstyle=\color{keyword}\bfseries,
  commentstyle=\color{comment}\itshape,
  stringstyle=\color{string},
  backgroundcolor=\color{codebg},
  frame=single,
  rulecolor=\color{codeframe},
  frameround=tttt,
  breaklines=true,
  showstringspaces=false,
  tabsize=2
}
\lstdefinestyle{Rstyle}{
  language=R,
  basicstyle=\ttfamily\small,
  keywordstyle=\color{keyword}\bfseries,
  commentstyle=\color{comment}\itshape,
  stringstyle=\color{string},
  backgroundcolor=\color{codebg},
  frame=single,
  rulecolor=\color{codeframe},
  frameround=tttt,
  breaklines=true,
  showstringspaces=false,
  tabsize=2
}
\lstset{style=pythonstyle}

% --- Common macros (configurable path) ---
\providecommand{\macrosPath}{../../preamble/macros.tex}
\input{\macrosPath}

% --- Exercise counter ---
\newcounter{excounter}
\renewcommand{\theexcounter}{\arabic{excounter}}

\newtcolorbox[use counter=excounter]{exercise}[1][]{
  enhanced, breakable,
  colback=boxEx!5, colframe=boxEx, fonttitle=\bfseries,
  title={Exercise~\theexcounter\ifstrempty{#1}{}{ -- #1}},
  arc=2pt, boxrule=0.6pt
}
\newtcolorbox{question}[1][]{
  enhanced, breakable,
  colback=boxQ!5, colframe=boxQ, fonttitle=\bfseries,
  title={Question\ifstrempty{#1}{}{ -- #1}},
  arc=2pt, boxrule=0.5pt
}
\newtcolorbox{solution}[1][]{
  enhanced, breakable,
  colback=boxC!5, colframe=boxC, fonttitle=\bfseries,
  title={Solution\ifstrempty{#1}{}{ -- #1}},
  arc=2pt, boxrule=0.5pt
}
\newtcolorbox{warning}[1][]{
  enhanced, breakable,
  colback=boxAtt!5, colframe=boxAtt, fonttitle=\bfseries,
  title={Warning\ifstrempty{#1}{}{ -- #1}},
  arc=2pt, boxrule=0.5pt
}

% --- Headers / footers ---
\pagestyle{fancy}
\fancyhf{}
\fancyhead[L]{\textsc{\moduleNom} -- \auteurNom}
\fancyhead[R]{\TPnumero}
\fancyfoot[C]{\thepage}
\fancyfoot[L]{\auteurInstitution}
\fancyfoot[R]{\auteurEmail}
\renewcommand{\headrulewidth}{0.4pt}
\renewcommand{\footrulewidth}{0.2pt}

% --- Placeholder ---
\providecommand{\TPnumero}{Lab}


\title{Project 3: Stationarity diagnostic on a real series\\
\large Decomposition $+$ ACF/PACF $+$ Ljung-Box}
\author{\auteurNom\\\auteurEmail\\\auteurInstitution}
\date{\anneeUniversitaire}

\begin{document}
\maketitle

\section*{Context}

This project builds on Project 2 (decomposition of \texttt{SouvenirSales.csv})
and asks you to assess the stationarity of the residual after
log-decomposition. You will produce a complete stationarity diagnostic
combining visual tools (ACF, PACF) and a formal test (Ljung--Box).

The point is to determine whether the residual is white noise (in
which case decomposition is sufficient) or whether it carries
exploitable autocorrelation that should be modelled with ARMA in
Chapter 5.

\section*{Deliverables}
\begin{itemize}
  \item Reproducible \texttt{project3\_<name>.py} with seed $42$.
  \item Report PDF of 4--6 pages including all figures and a final
        recommendation.
  \item Optional 5-min oral defence.
\end{itemize}

% =========================================================
\begin{exercise}[Setup and decomposition]
\textbf{(1)} Load \texttt{data/SouvenirSales.csv}, take logs.

\textbf{(2)} On $\log Y_t$, fit a model
$\log Y_t = \beta_0 + \beta_1\,t + \sum_{j=1}^{2}[a_j \cos(j\omega t)
+ b_j \sin(j\omega t)] + \varepsilon_t$ with $\omega = 2\pi / 12$ by OLS.

\textbf{(3)} Compute the residual $\widehat\varepsilon_t$. Plot it
versus $t$. Comment on its visual stationarity (mean stable around
$0$? variance constant?).
\end{exercise}

% =========================================================
\begin{exercise}[Sample autocorrelation]
\textbf{(1)} Plot the sample ACF of $\widehat\varepsilon_t$ up to lag
24, with Bartlett bands.

\textbf{(2)} Plot the sample PACF up to lag 24.

\textbf{(3)} From the ACF/PACF patterns, propose a tentative ARMA
order $(p, q)$. Justify your choice using the signature table of
Chapter 3.
\end{exercise}

% =========================================================
\begin{exercise}[Ljung--Box test]
\textbf{(1)} Run the Ljung--Box test on $\widehat\varepsilon_t$ at
lags $m \in \{6, 12, 18, 24\}$.

\textbf{(2)} Build a small table with columns: $m$, $Q_{\text{LB}}$,
critical value $\chi^2_{m, 0.95}$, $p$-value, decision.

\textbf{(3)} Conclude: is $\widehat\varepsilon_t$ white noise at the
$5\%$ level? Why does the answer typically differ between $m = 6$
and $m = 24$?
\end{exercise}

% =========================================================
\begin{exercise}[Comparison: pre-decomposition vs.\ post-decomposition]
\textbf{(1)} Run the same Ljung--Box test on the raw $\log Y_t$ (no
decomposition) at the same lags. What happens?

\textbf{(2)} Compute the variance ratio
$\Var[\widehat\varepsilon] / \Var[\log Y]$. What does it tell you
about the explanatory power of the decomposition?

\textbf{(3)} Discuss in 5 lines: even if the residual fails the
Ljung--Box test, the decomposition is still useful — why?
\end{exercise}

% =========================================================
\begin{exercise}[Bonus: AR(1) modelling of the residual]
If the Ljung--Box test rejects white noise, fit an AR(1) on
$\widehat\varepsilon_t$:
\[
  \widehat\varepsilon_t = \phi\, \widehat\varepsilon_{t-1} + \eta_t.
\]
\textbf{(1)} Estimate $\widehat\phi$ by OLS.

\textbf{(2)} Plot the residuals $\widehat\eta_t$ and re-run the
Ljung--Box test.

\textbf{(3)} Has the AR(1) absorbed the remaining autocorrelation?
\end{exercise}

% =========================================================
\section*{Marking grid (out of 100)}

\begin{center}
\begin{tabular}{p{8cm}rl}
\toprule
\textbf{Criterion} & \textbf{Points} & \\
\midrule
OLS decomposition + residual plot                & 15 \\
ACF / PACF visualisation + Box-Jenkins reading   & 25 \\
Ljung-Box table + interpretation                 & 25 \\
Comparison pre/post + variance reduction         & 15 \\
Bonus AR(1) fit + residual diagnostics           & 10 \\
Reproducibility + report quality                 & 10 \\
\midrule
\textbf{Total}                                   & \textbf{100} \\
\bottomrule
\end{tabular}
\end{center}

\bigskip
\hfill\auteurNom\par
\hfill\auteurEmail\par

\end{document}
